% META: {"id": "1SPE-PROBCOND-EX-018", "chapitre": "1SPE-PROBA-COND", "type_objet": "exercice", "capacites": ["1SPE-PROBA-COND-C2"], "capacites_codes": ["C2"], "methodes": ["M2"], "parcours": 2, "competences": ["calculer", "raisonner"], "duree_min": 15, "mode_creation": "ex_nihilo", "sources_inspiration": [], "parametres_sympy": {}, "fichier_tex": "chapitres/1SPE-PROBA-COND/exercices/1SPE-PROBCOND-EX-018.tex", "corrige_tex": "chapitres/1SPE-PROBA-COND/corriges/1SPE-PROBCOND-CO-018.tex", "status": "generated"}
% BEGIN-VERIFY
% from sympy import Rational
% # 3 urnes. Choix equi. U1: 2R,1B. U2: 1R,2B. U3: 3R,0B.
% P_U = Rational(1, 3)
% P_R_U1 = Rational(2, 3)
% P_R_U2 = Rational(1, 3)
% P_R_U3 = Rational(1, 1)
% P_R = P_U*(P_R_U1 + P_R_U2 + P_R_U3)
% assert P_R == Rational(1,3)*(Rational(2,3) + Rational(1,3) + 1)
% assert P_R == Rational(1,3) * 2
% assert P_R == Rational(2, 3)
% END-VERIFY
\begin{exercice}{1SPE-PROBCOND-EX-018}{2}{15}
On dispose de trois urnes identiques. L'urne $U_1$ contient $2$ boules rouges et $1$ bleue. L'urne $U_2$ contient $1$ rouge et $2$ bleues. L'urne $U_3$ contient $3$ rouges. On choisit une urne au hasard (equiprobabilite) puis on tire une boule.

\begin{enumerate}
  \item Construire l'arbre pondere.
  \item Calculer la probabilite de tirer une boule rouge.
\end{enumerate}
\end{exercice}
